(1) 时间独立薛定谔方程为 \(-\hbar^2\psi''/(2m)+V\psi=E\psi\)。在势垒区 \(V_0>E\)，波数平方为 \(2m(E-V_0)/\hbar^2<0\)，因此解为增长与衰减指数的线性组合，定义 \(\kappa=\sqrt{2m(V_0-E)}/\hbar>0\)，振幅典型衰减为 \(e^{-\kappa x}\)。
(2) 一般WKB禁阻区穿透率的指数为 \(\exp[-(2/\hbar)\int\sqrt{2m(V(x)-E)}dx]\)。矩形势垒宽度 \(a\) 内根号为常数，得到 \(T_{\rm WKB}\sim e^{-2\kappa a}\)。这里仅给出指数控制项，不保证前因子为一，尤其矩形跃变边界并非缓变势。
(3) 对矩形势垒两边的连续波函数和一阶导数匹配，得准确系数 \(T=[1+V_0^2\sinh^2(\kappa a)/(4E(V_0-E))]^{-1}\)。当 \(\kappa a\gg1\) 时 \(\sinh^2(\kappa a)\sim e^{2\kappa a}/4\)，故 \(T\sim16E(V_0-E)e^{-2\kappa a}/V_0^2\)，指数与WKB一致，而前因子不同于一。
(4) 当 \(a\to0\) 时 \(\sinh(\kappa a)\to0\)，精确式给 \(T\to1\)；固定 \(a\) 取 \(V_0\to E^+\)，应同时展开 \(\kappa\)，而非仅把WKB指数视为完整结果。增大 \(m\) 会增大 \(\kappa\)，抑制隧穿；缓变势的WKB公式在急剧跳变的边界只能用于主指数的启发式估计。
